\setcounter{section}{16}






  \zwischenueberschrift{Irreduzible Filter}

\bild{ \begin{center}
\includegraphics[width=5.5cm]{ \bildeinlesungjpg  {Kaffeefilter} {jpg } }
\end{center}
\bildtext {Einen Filter kann man mit dem, was in ihm hängen bleibt, identifizieren.} }
\bildlizenz { Kaffeefilter.jpg  } {Elke Wetzig} {Elya} {Commons} {CC-BY-SA-3.0} {}




Wir haben in der letzten Vorlesung gesehen, dass zu einem Punkt $P$ in einem $K$-Spektrum
\mathl{K\!-\!\operatorname{Spek}\, \left(  R  \right)}{} der Umgebungsfilter gehört und dass der Halm in diesem Filter gleich der Lokalisierung von $R$ an dem zugehörigen maximalen Ideal

\mavergleichskette
{\vergleichskette
{ {\mathfrak m}_P
}
{ \subseteq }{ R
}
{  }{
}
{    }{
}
{   }{
}
}
{}{}{}
ist. Ebenfalls haben wir gesehen, dass bei integrem $R$ der Halm über alle nichtleeren offenen Mengen den Quotientenkörper von $R$ liefert, der wiederum die Lokalisierung am Nullideal ist. Dieser Zusammenhang wird mit dem Begriff des irreduziblen Filters verallgemeinert.


\inputdefinition
{}
{





Ein \definitionsverweis {topologischer Filter}{}{}
$F$ heißt \definitionswort {irreduzibel}{,} wenn

\mavergleichskette
{\vergleichskette
{  \emptyset
}
{ \notin }{ F
}
{  }{
}
{    }{
}
{   }{
}
}
{}{}{}
und folgendes gilt: Sind $U,V$ zwei
\definitionsverweis {offene Mengen}{}{}
mit

\mavergleichskette
{\vergleichskette
{  U \cup V
}
{ \in }{ F
}
{  }{
}
{    }{
}
{   }{
}
}
{}{}{,}
so ist
\mathkor {} {U \in F} {oder} {V \in F} {.}





}

Für Zariski-Filter
\zusatzklammer {also topologische Filter in der Zariski-Topologie} {} {}
gilt folgender Zusammenhang.



\inputfaktbeweis
{K-Spektrum/Irreduzible Filter, Primideale, irreduzible Teilmengen/Fakt}
{Satz}
{}
{



\faktsituation {}
\faktvoraussetzung {Es sei $K$ ein
\definitionsverweis {algebraisch abgeschlossener Körper}{}{}
und sei $R$ eine kommutative
$K$-\definitionsverweis {Algebra von endlichem Typ}{}{}
mit
$K$-\definitionsverweis {Spektrum}{}{}

\mavergleichskette
{\vergleichskette
{ X
}
{ = }{ K\!-\!\operatorname{Spek}\, \left(  R  \right)
}
{  }{
}
{    }{
}
{   }{
}
}
{}{}{.}}
\faktfolgerung {Dann entsprechen sich folgende Objekte.
\aufzaehlungdrei{ \definitionsverweis {Primideale}{}{}
in $R$.
}{ \definitionsverweis {Irreduzible}{}{}
\definitionsverweis {abgeschlossene Teilmengen}{}{}
von $X$.
}{ \definitionsverweis {Irreduzible Filter}{}{}
in $X$.
}}
\faktzusatz {Dabei entspricht der irreduziblen abgeschlossenen Teilmenge

\mavergleichskette
{\vergleichskette
{ Y
}
{ \subseteq }{ X
}
{  }{
}
{    }{
}
{   }{
}
}
{}{}{}
der Filter

\mavergleichskettedisp
{\vergleichskette
{ F(Y)
}
{ =} { \left\{  U \subseteq X   \mid   U \cap Y \neq \emptyset        \right\}
}
{ } {
}
{ } {
}
{ } {
}
}
{}{}{.}
Der
\definitionsverweis {Halm}{}{}
der Strukturgarbe an diesem Filter ist die Lokalisierung $R_{\mathfrak p}$, wobei ${\mathfrak p}$ das zugehörige Primideal bezeichnet.}
\faktzusatz {}





}
{





Die Korrespondenz zwischen Primidealen und abgeschlossenen irreduziblen Teilmengen
\zusatzklammer {zu einem Primideal ${\mathfrak p}$ gehört die irreduzible abgeschlossene Teilmenge \mathlk{V({\mathfrak p})}{}} {} {}
ist bekannt
\zusatzklammer {siehe
Lemma 4.3

und

Proposition 11.7
} {} {.}
Die angegebene Konstruktion zu einer irreduziblen abgeschlossenen Menge $Y$ liefert in der Tat einen irreduziblen Filter. Dabei ist die Irreduzibilität trivial, zu zeigen ist lediglich die Durchschnittseigenschaft eines Filters. Es seien

\mavergleichskette
{\vergleichskette
{ U, V
}
{ \in }{ F
}
{ = }{ F(Y)
}
{    }{
}
{   }{
}
}
{}{}{,}
sodass also
die Durchschnitte
\mathl{Y \cap U}{} und
\mathl{Y \cap V}{} nicht leer sind. Dann ist aber wegen der Irreduzibilität von $Y$ auch der Durchschnitt

\mavergleichskettedisp
{\vergleichskette
{ (Y \cap U) \cap (Y \cap V)
}
{ =} { Y \cap (U \cap V)
}
{ } {
}
{ } {
}
{ } {
}
}
{}{}{}
nicht leer und daher ist

\mavergleichskette
{\vergleichskette
{ U \cap V
}
{ \in }{ F
}
{  }{
}
{    }{
}
{   }{
}
}
{}{}{.}

Es sei nun $F$ irgendein irreduzibler topologischer Filter. Wir behaupten, dass das Komplement von

\mavergleichskettedisp
{\vergleichskette
{S
}
{ =} { \left\{  f \in R   \mid   D(f)\in F        \right\}
}
{ } {
}
{ } {
}
{ } {
}
}
{}{}{}
ein Primideal ist. Es ist sofort ein saturiertes multiplikatives System in $R$. Es bleibt zu zeigen, dass das Komplement additiv abgeschlossen ist. Es seien dazu

\mavergleichskette
{\vergleichskette
{ h,g
}
{ \in }{ R
}
{  }{
}
{    }{
}
{   }{
}
}
{}{}{}
mit

\mavergleichskette
{\vergleichskette
{ g+h
}
{ \in }{ S
}
{  }{
}
{    }{
}
{   }{
}
}
{}{}{,}
also

\mavergleichskette
{\vergleichskette
{ D(g+h)
}
{ \in }{ F
}
{  }{
}
{    }{
}
{   }{
}
}
{}{}{.}
Dann gehört erst recht

\mavergleichskettedisp
{\vergleichskette
{ D(g) \cup D(h)
}
{ =} { D(g,h)
}
{ \supseteq} { D(g+h)
}
{ } {
}
{ } {
}
}
{}{}{}
zu $F$ und wegen der Irreduzibilität von $F$ ist

\mavergleichskette
{\vergleichskette
{ D(g)
}
{ \in }{ F
}
{  }{
}
{    }{
}
{   }{
}
}
{}{}{}
oder

\mavergleichskette
{\vergleichskette
{ D(h)
}
{ \in }{ F
}
{  }{
}
{    }{
}
{   }{
}
}
{}{}{,}
woraus sich

\mavergleichskette
{\vergleichskette
{ g
}
{ \in }{ S
}
{  }{
}
{    }{
}
{   }{
}
}
{}{}{}
oder

\mavergleichskette
{\vergleichskette
{ h
}
{ \in }{ S
}
{  }{
}
{    }{
}
{   }{
}
}
{}{}{}
ergibt.

Diese drei Zuordnungen hintereinandergenommen führen dabei immer wieder zum Ausgangsobjekt zurück. Dazu muss man lediglich beachten, dass ein irreduzibler Zariski-Filter durch offene Mengen der Form
\mathl{D(f)}{} erzeugt wird, siehe

Aufgabe 16.1
.
Der Zusatz ist ein Spezialfall von

Aufgabe 15.27
.




}



Den zu einer irreduziblen abgeschlossenen Menge gehörenden Filter nennen wir auch den zugehörigen \stichwort {generischen Filter} {} zu $Y$ und den Halm davon den \stichwort {generischen Halm} {} zu $Y$. Ein Spezialfall der Korrespondenz von

Satz 16.2

ist die Beziehung zwischen minimalen Primidealen, irreduziblen Komponenten und Ultrafiltern. Auf der anderen Seite hat man die Korrespondenz zwischen maximalen Idealen, Punkten und Umgebungsfiltern.





  \zwischenueberschrift{Morphismen zwischen Varietäten}


\inputdefinition
{}
{





Es seien $X$ und $Y$
\definitionsverweis {quasiaffine Varietäten}{}{}
und sei
\maabbdisp {{\psi}} { Y } { X
} {}
eine stetige Abbildung. Dann nennt man ${\psi}$ einen \definitionswort {Morphismus}{}
\zusatzklammer {von quasiaffinen Varietäten} {} {,}
wenn für jede offene Teilmenge

\mavergleichskette
{\vergleichskette
{ U
}
{ \subseteq }{ X
}
{  }{
}
{    }{
}
{   }{
}
}
{}{}{}
und jede
\definitionsverweis {algebraische Funktion}{}{}

\mavergleichskette
{\vergleichskette
{ f
}
{ \in }{ \Gamma ( U , {\mathcal O} )
}
{  }{
}
{    }{
}
{   }{
}
}
{}{}{}
gilt, dass die zusammengesetzte Funktion

\mathdisp {f \circ {\psi} : {\psi}^{-1}(U)  \longrightarrow U \stackrel{f}{\longrightarrow} {\mathbb A}^{1}_{K}} {  }

zu
\mathl{\Gamma ( {\psi}^{-1}(U), {\mathcal O} )}{} gehört.





}

\inputbemerkung
{}
{





Ein Morphismus
\maabbdisp {{\psi}} { Y } { X
} {}
induziert nach Definition zu jeder offenen Teilmenge

\mavergleichskette
{\vergleichskette
{ U
}
{ \subseteq }{ X
}
{  }{
}
{    }{
}
{   }{
}
}
{}{}{}
einen
\definitionsverweis {Ringhomomorphismus}{}{}
\maabbdisp {\tilde{\psi}} { \Gamma ( U , {\mathcal O} ) } { \Gamma ( {\psi}^{-1}(U), {\mathcal O} )
} {,}
Insbesondere gibt es einen \stichwort {globalen Ringhomomorphismus} {}
\maabbdisp {\tilde{\psi}} { \Gamma ( X , {\mathcal O} ) } { \Gamma ( Y, {\mathcal O} )
} {.}
Sind

\mavergleichskette
{\vergleichskette
{ U_1
}
{ \subseteq }{ U_2
}
{  }{
}
{    }{
}
{   }{
}
}
{}{}{}
offene Teilmengen in $Y$, so liegt ein kommutatives Diagramm von stetigen Abbildungen vor
\zusatzklammer {wobei die senkrechten Pfeile offene Inklusionen sind} {} {}

\mathdisp {\begin{matrix}



{\psi}^{-1}(U_1)  & \stackrel{  }{\longrightarrow} &  U_1
&  \\ \downarrow & & \downarrow  & \\  {\psi}^{-1}(U_2)  & \stackrel{  }{\longrightarrow} &  U_2
& \! \! \! \!\!\!\!\! ,  \\ \end{matrix}} {  }




das wiederum zu dem kommutativen Diagramm

\mathdisp {\begin{matrix}

\Gamma ( {\psi}^{-1}(U_1), {\mathcal O} )  & \stackrel{}{\longleftarrow} &  \Gamma ( U_1 , {\mathcal O} )



&  \\



\uparrow &  & \uparrow & \\



\Gamma ( {\psi}^{-1}(U_2), {\mathcal O} )  & \stackrel{}{\longleftarrow} &  \Gamma ( U_2 , {\mathcal O} )
& \!\!\!\!\!    \\



\end{matrix}} {  }

von Ringhomomorphismen führt.





}

Wir fassen einige einfache Eigenschaften von Morphismen zusammen.



\inputfaktbeweistrivial
{K-Spektrum/Morphismus/Verknüpfung/Offene Einbettung/Fakt}
{Proposition}
{}
{



\faktsituation {}
\faktvoraussetzung {Es sei $K$ ein algebraisch abgeschlossener Körper, und seien
\mathl{U,X,Y,Z}{} quasiaffine Varietäten.}
\faktuebergang {Dann gelten folgende Aussagen.}
\faktfolgerung {\aufzaehlungzwei {Eine offene Einbettung

\mavergleichskette
{\vergleichskette
{ U
}
{ \subseteq }{ X
}
{  }{
}
{    }{
}
{   }{
}
}
{}{}{}
ist ein
\definitionsverweis {Morphismus}{}{.}
} {Sind
\mathl{{\theta}:Z \rightarrow Y}{} und
\mathl{{\psi}:Y \rightarrow X}{} Morphismen, so ist auch die Verknüpfung
\mathl{{\psi} \circ {\theta}}{} ein Morphismus.
}}
\faktzusatz {}
\faktzusatz {}





}







Wichtiger sind die folgenden Eigenschaften.



\inputfaktbeweis
{K-Spektrum/Ringhomomorphismus induziert Morphismus/Fakt}
{Satz}
{}
{



\faktsituation {}
\faktvoraussetzung {Es sei $K$ ein
\definitionsverweis {algebraisch abgeschlossener Körper}{}{}
und seien $R$ und $S$ kommutative
$K$-\definitionsverweis {Algebren vom endlichen Typ}{}{}
mit zugehörigen
$K$-\definitionsverweis {Spektren}{}{}

\mathl{X= K\!-\!\operatorname{Spek}\, \left(  R  \right)}{} und
\mathl{Y= K\!-\!\operatorname{Spek}\, \left(  S  \right)}{.}}
\faktfolgerung {Dann ist die durch einen
$K$-\definitionsverweis {Algebrahomomorphismus}{}{}
\maabb {\varphi} { R } { S
} {}
induzierte
\definitionsverweis {Spektrumsabbildung}{}{}
\maabbdisp {\varphi^*} { Y} { X
} {}
ein
\definitionsverweis {Morphismus}{}{.}}
\faktzusatz {}
\faktzusatz {}





}
{





Wir wissen bereits nach

Satz 12.7
,
dass
\maabbdisp {} { Y = K\!-\!\operatorname{Spek}\, \left(  S  \right) } { X = K\!-\!\operatorname{Spek}\, \left(  R  \right)
} {}
eine stetige Abbildung ist. Es sei

\mavergleichskette
{\vergleichskette
{ U
}
{ \subseteq }{ X
}
{  }{
}
{    }{
}
{   }{
}
}
{}{}{}
eine offene Teilmenge und

\mavergleichskette
{\vergleichskette
{ V
}
{ = }{ (\varphi^*)^{-1}(U)
}
{  }{
}
{    }{
}
{   }{
}
}
{}{}{}
das Urbild. Es sei
\maabb {f} { U } { K
} {}
eine
\definitionsverweis {algebraische Funktion
}{}{}
mit der Hintereinanderschaltung
\maabb {f \circ \varphi^*} { V } { K
} {.}
Wir müssen zeigen, dass diese Abbildung ebenfalls algebraisch ist. Es sei dazu

\mavergleichskette
{\vergleichskette
{ P
}
{ \in }{ V
}
{  }{
}
{    }{
}
{   }{
}
}
{}{}{}
ein Punkt mit dem Bildpunkt

\mavergleichskette
{\vergleichskette
{ Q
}
{ = }{ \varphi^*(P)
}
{  }{
}
{    }{
}
{   }{
}
}
{}{}{.}
Es sei

\mavergleichskette
{\vergleichskette
{ Q
}
{ \in }{ D(H)
}
{ \subseteq }{ U
}
{    }{
}
{   }{
}
}
{}{}{}
und

\mavergleichskette
{\vergleichskette
{ f
}
{ = }{ G/H
}
{  }{
}
{    }{
}
{   }{
}
}
{}{}{}
auf $D(H)$ mit

\mavergleichskette
{\vergleichskette
{ G, H
}
{ \in }{ R
}
{  }{
}
{    }{
}
{   }{
}
}
{}{}{.}
Es ist

\mavergleichskettedisp
{\vergleichskette
{ P
}
{ \in} { (\varphi^*)^{-1}( D(H))
}
{ =} { D( \varphi (H))
}
{ } {
}
{ } {
}
}
{}{}{}
gemäß

Satz 12.7
.
Wir behaupten, dass auf
\mathl{D( \varphi (H))}{} die Gleichheit

\mavergleichskette
{\vergleichskette
{ f \circ \varphi^*
}
{ = }{ { \frac{   \varphi(G)  }{  \varphi(H)  } }
}
{  }{
}
{    }{
}
{   }{
}
}
{}{}{}
gilt. Dies folgt für

\mavergleichskette
{\vergleichskette
{ \tilde{P}
}
{ \in }{ D( \varphi (H))
}
{  }{
}
{    }{
}
{   }{
}
}
{}{}{}
aus

\mavergleichskettedisp
{\vergleichskette
{ f \circ \varphi^*   (\tilde{P})
}
{ =} { f( \varphi^*(\tilde{P}) )
}
{ =} { \frac{ G( \varphi^*( \tilde{P} ) )}{ H( \varphi^*( \tilde{P}) ) }
}
{ =} { \frac{ ( \varphi(G)) (\tilde{P})}{( \varphi(H)) ( \tilde{P} )}
}
{ } {
}
}
{}{}{.}




}




\inputbemerkung
{}
{





In der Situation von

Satz 16.6

ist der zu

\mavergleichskette
{\vergleichskette
{ U
}
{ = }{ D(f)
}
{  }{
}
{    }{
}
{   }{
}
}
{}{}{}
gehörende
\definitionsverweis {Ringhomomorphismus}{}{}
die natürliche Abbildung
\maabbdisp {} { \Gamma (D(f), {\mathcal O} ) \cong R_f } { \Gamma ((\varphi^*)^{-1}(D(f)), {\mathcal O} ) = \Gamma (D(\varphi(f)), {\mathcal O} ) =  S_{\varphi(f)}
} {.}





}





\inputfaktbeweis
{Quasiaffine Varietät/Globale algebraische Funktion/Ist Morphismus nach affiner Geraden/Fakt}
{Lemma}
{}
{



\faktsituation {}
\faktvoraussetzung {Es sei $U$ eine
\definitionsverweis {quasiaffine Varietät}{}{}
über einem
\definitionsverweis {algebraisch abgeschlossenen Körper}{}{}
$K$ und sei

\mavergleichskette
{\vergleichskette
{ f
}
{ \in }{ \Gamma (U, {\mathcal O} )
}
{  }{
}
{    }{
}
{   }{
}
}
{}{}{}
eine
\definitionsverweis {algebraische Funktion}{}{.}}
\faktfolgerung {Dann definiert $f$ einen
\definitionsverweis {Morphismus}{}{}
\maabbdisp {f} { U } { {\mathbb A}^{1}_{ K }  \cong K
} {.}}
\faktzusatz {}
\faktzusatz {}





}
{





Eine stetige Abbildung liegt aufgrund von

Aufgabe 14.4

vor. Es sei

\mavergleichskette
{\vergleichskette
{ U
}
{ \subseteq }{ K\!-\!\operatorname{Spek}\, \left(  R  \right)
}
{  }{
}
{    }{
}
{   }{
}
}
{}{}{.}
Es sei

\mavergleichskette
{\vergleichskette
{ V
}
{ = }{ D(s)
}
{ \subseteq }{ {\mathbb A}^{1}_{K}
}
{    }{
}
{   }{
}
}
{}{}{}
eine offene Teilmenge und

\mavergleichskette
{\vergleichskette
{ W
}
{ = }{ f^{-1}(V)
}
{ \subseteq }{ U
}
{    }{
}
{   }{
}
}
{}{}{}
das Urbild davon. Es sei

\mavergleichskettedisp
{\vergleichskette
{ q
}
{ =} { { \frac{   r  }{  s^n  } }
}
{ \in} { \Gamma (V, {\mathcal O} )
}
{ =} { K[T]_{ s }
}
{ } {
}
}
{}{}{}
eine algebraische Funktion auf $V$. Wir müssen zeigen, dass die Verknüpfung
\mathl{q \circ f}{} eine algebraische Funktion auf $W$ ist. Es sei dazu

\mavergleichskette
{\vergleichskette
{ P
}
{ \in }{ W
}
{  }{
}
{    }{
}
{   }{
}
}
{}{}{}
und sei

\mavergleichskette
{\vergleichskette
{ f
}
{ = }{ G/H
}
{  }{
}
{    }{
}
{   }{
}
}
{}{}{}
eine Beschreibung der nach Voraussetzung algebraischen Funktion $f$ in der Umgebung

\mavergleichskette
{\vergleichskette
{ D(H)
}
{ \ni }{ P
}
{  }{
}
{    }{
}
{   }{
}
}
{}{}{.}
Dann ist

\mavergleichskettedisp
{\vergleichskette
{ ( q \circ f ) (P)
}
{ =} { q (f(P))
}
{ =} { \frac{r}{s^n} \left(  \frac{G}{H} (P)  \right)
}
{ =} { \frac{ r( G(P)/H(P) )}{(s(G(P)/H(P)))^n}
}
{ } {
}
}
{}{}{.}
Dabei ist der Nenner
\mathl{(s(G(P)/H(P)))^n}{} nicht $0$, da

\mavergleichskette
{\vergleichskette
{ f(P)
}
{ \in }{  D(s)
}
{  }{
}
{    }{
}
{   }{
}
}
{}{}{}
ist, sodass dies eine rationale Darstellung ist.




}






\inputfaktbeweis
{Quasiaffine Varietät/Morphismus nach affiner Geraden/Fakt}
{Satz}
{}
{



\faktsituation {}
\faktvoraussetzung {Es sei $U$ eine quasiaffine Varietät über einem
\definitionsverweis {algebraisch abgeschlossenen Körper}{}{}
$K$, und zwar sei

\mavergleichskette
{\vergleichskette
{ U
}
{ \subseteq }{ K\!-\!\operatorname{Spek}\, \left(  R  \right)
}
{  }{
}
{    }{
}
{   }{
}
}
{}{}{,}
wobei $R$ eine kommutative $K$-Algebra von endlichem Typ sei.}
\faktfolgerung {Dann gibt es eine natürliche Bijektion
\maabbeledisp {} { \operatorname{Mor} \,( U ,  {\mathbb A}^{1}_{K} ) } { \Gamma (U, {\mathcal O} )
} { {\psi}  } { \tilde{\psi} (T)
} {,}
wobei $T$ die Variable in

\mavergleichskette
{\vergleichskette
{ K[T]
}
{ = }{ \Gamma ( {\mathbb A}^{1}_{K} , {\mathcal O} )
}
{  }{
}
{    }{
}
{   }{
}
}
{}{}{}
bezeichnet.}
\faktzusatz {Insbesondere sind Morphismen von $U$ in die affine Gerade durch den globalen Ringhomomorphismus
\maabbdisp {\tilde{\psi}} { \Gamma ( {\mathbb A}^{1}_{K}, {\mathcal O} ) = K[T] } { \Gamma (U, {\mathcal O} )
} {}
eindeutig bestimmt.}
\faktzusatz {}





}
{





Die Abbildung ist wohldefiniert und surjektiv. Ist nämlich eine globale algebraische Funktion

\mavergleichskette
{\vergleichskette
{ f
}
{ \in }{ \Gamma (U, {\mathcal O} )
}
{  }{
}
{    }{
}
{   }{
}
}
{}{}{}
gegeben, so ist zunächst
\maabb {f} { U } { K
} {.}
Die Variable $T$, die auf

\mavergleichskette
{\vergleichskette
{ K
}
{ = }{ {\mathbb A}^{1}_{K}
}
{  }{
}
{    }{
}
{   }{
}
}
{}{}{}
der identischen Abbildung entspricht, wird unter
\zusatzklammer {der Verknüpfung mit} {} {}
$f$ auf das Element

\mavergleichskette
{\vergleichskette
{ f
}
{ \in }{ \Gamma (U, {\mathcal O} )
}
{  }{
}
{    }{
}
{   }{
}
}
{}{}{}
abgebildet. Nach

Lemma 16.8

ist $f$ ein Morphismus.

Die Injektivität ergibt sich, da sowohl der Morphismus als auch die algebraische Funktion durch die zugrunde liegende stetige Abbildung eindeutig festgelegt sind.




}






\inputfaktbeweis
{Quasiaffine Varietät/Morphismus nach affiner Varietät/Fakt}
{Satz}
{}
{



\faktsituation {}
\faktvoraussetzung {Es sei $U$ eine quasiaffine Varietät, und zwar sei

\mavergleichskette
{\vergleichskette
{U
}
{ \subseteq }{ K\!-\!\operatorname{Spek}\, \left(  R  \right)
}
{  }{
}
{    }{
}
{   }{
}
}
{}{}{,}
wobei $R$ eine kommutative $K$-Algebra von endlichem Typ über einem algebraisch abgeschlossener Körper $K$ sei. Es sei $S$ eine weitere kommutative $K$-Algebra von endlichem Typ.}
\faktfolgerung {Dann gibt es eine natürliche Bijektion
\maabbeledisp {} { \operatorname{Mor} \,( U ,  K\!-\!\operatorname{Spek}\, \left(  S  \right) ) } { \operatorname{Hom}^{\rm alg}_{ K } \, \left(   S , \Gamma (U, {\mathcal O} )   \right)
} { {\psi} } { \tilde{\psi}
} {,}
wobei $\tilde{\psi}$ den zu ${\psi}$ gehörigen globalen Ringhomomorphismus bezeichnet.}
\faktzusatz {}
\faktzusatz {}





}
{





Die Abbildung ist wohldefiniert. Aus

Satz 16.9

folgt, dass die Aussage für

\mavergleichskette
{\vergleichskette
{ S
}
{ = }{ K[T]
}
{  }{
}
{    }{
}
{   }{
}
}
{}{}{}
richtig ist. Daraus ergibt sich, dass die Aussage für jeden Polynomring
\mathl{K[T_1 , \ldots , T_n]}{} richtig ist, da ein Morphismus nach ${ {\mathbb A}_{  K }^{  n   } }$ durch seine Komponenten und ein $K$-Algebrahomomorphismus durch die Einsetzungen für $T_i$ gegeben ist. Es sei nun

\mavergleichskette
{\vergleichskette
{ S
}
{ = }{ K [ T_1, \ldots ,  T_{ n }]/{\mathfrak a}
}
{  }{
}
{    }{
}
{   }{
}
}
{}{}{}
und

\mavergleichskettedisp
{\vergleichskette
{ K\!-\!\operatorname{Spek}\, \left(  S  \right)
}
{ \cong} { V( {\mathfrak a} )
}
{ =} { V
}
{ \subseteq} { { {\mathbb A}_{  K }^{  n   } }
}
{ } {
}
}
{}{}{.}
Zu einem Morphismus
\maabb {} { U } { K\!-\!\operatorname{Spek}\, \left(  S  \right)
} {}
ist die Verknüpfung mit der abgeschlossenen Einbettung in den affinen Raum ebenfalls ein Morphismus. D.h. es liegt ein kommutatives Diagramm

\mathdisp {\begin{matrix}
\operatorname{Mor} \,( U ,  K\!-\!\operatorname{Spek}\, \left(  S  \right))  & \stackrel{  }{\longrightarrow} &  \operatorname{Hom}^{\rm alg}_{ K } \, \left(   S , \Gamma (U, {\mathcal O} )  \right)



&  \\ \downarrow & & \downarrow  & \\  \operatorname{Mor} \,( U ,  { {\mathbb A}_{  K }^{  n   } })  & \stackrel{  }{\longrightarrow} &  \operatorname{Hom}^{\rm alg}_{ K } \, \left(   K [ T_1, \ldots ,  T_{ n }] , \Gamma (U, {\mathcal O} )  \right)
& \! \! \! \!\!\!\!\! ,  \\ \end{matrix}} {  }




vor, wobei die untere Abbildung bereits als Bijektion nachgewiesen wurde. Die vertikalen Abbildungen sind injektiv. Wir müssen daher zeigen, dass die untere Abbildung die oberen Teilmengen ineinander überführt.

Ein Morphismus
\maabb {} { U } { { {\mathbb A}_{  K }^{  n   } }
} {,}
der
\zusatzklammer {als Abbildung} {} {}
durch $V$ faktorisiert, ist auch ein Morphismus nach $V$. Die Morphismuseigenschaft ist nur für offene Mengen der Form $D(H)$ zu überprüfen,

\mavergleichskette
{\vergleichskette
{ H
}
{ \in }{ S
}
{  }{
}
{    }{
}
{   }{
}
}
{}{}{.}
Es sei

\mavergleichskette
{\vergleichskette
{   \tilde{H}
}
{ \in }{ K [ T_1, \ldots ,  T_{ n }]
}
{  }{
}
{    }{
}
{   }{
}
}
{}{}{}
ein Repräsentant für $H$. Dann ist
\maabb {} { K [ T_1, \ldots ,  T_{ n }]_{\tilde{H} }} { S_H
} {}
surjektiv und damit wird jedes Element aus $S_H$ auf eine algebraische Funktion abgebildet.

Auf der rechten Seite des Diagramms gehört ein Algebrahomomorphismus genau dann zur oberen Menge, wenn ${\mathfrak a}$ zum Kern gehört. Damit folgt die Aussage aus

Aufgabe 16.8
.




}


\bild{ \begin{center}
\includegraphics[width=5.5cm]{ \bildeinlesungpng {Cone_intersects_line} {png} }
\end{center}
\bildtext {Nicht jede außerhalb der Gerade auf dem Kegel definierte Funktion lässt sich auf den affinen Raum ohne die Gerade fortsetzen.} }
\bildlizenz { Cone intersects line.png } {} {Pmidden} {Commons} {PD} {}





\inputbeispiel{}
{





Wir betrachten den Standardkegel, der als abgeschlossene Teilmenge

\mavergleichskettedisp
{\vergleichskette
{V
}
{ =} { V(X^2+Y^2-Z^2)
}
{ \subseteq} { { {\mathbb A}_{ K }^{  3   } }
}
{ } {
}
{ } {
}
}
{}{}{}
gegeben sei. Es sei

\mavergleichskette
{\vergleichskette
{U
}
{ = }{ D(X,Z-Y)
}
{ \subseteq }{ { {\mathbb A}_{ K }^{  3   } }
}
{    }{
}
{   }{
}
}
{}{}{}
die offene Teilmenge mit dem Durchschnitt

\mavergleichskette
{\vergleichskette
{ V \cap U
}
{ = }{ D(X,Z-Y)
}
{  }{
}
{    }{
}
{   }{
}
}
{}{}{}
\zusatzklammer {in $V$} {} {,}
der eine offene Menge in $V$ ist. Wir behaupten, dass der zugehörige Ringhomomorphismus
\maabbdisp {} { \Gamma (U, {\mathcal O} ) } { \Gamma ( U \cap V , {\mathcal O} )
} {}
nicht surjektiv ist. Das liegt daran, dass links einfach der Polynomring in drei Variablen steht
\zusatzklammer {vergleiche

Aufgabe 14.24
} {} {.}
Dagegen ergibt sich aus der Gleichung

\mavergleichskettedisp
{\vergleichskette
{ X^2
}
{ =} { Z^2-Y^2
}
{ =} { (Z-Y)(Z+Y)
}
{ } {
}
{ } {
}
}
{}{}{,}
dass es auf
\mathl{U \cap V}{} die algebraische Funktion

\mavergleichskettedisp
{\vergleichskette
{ \frac{X}{Z-Y}
}
{ =} { \frac{Z+Y}{X}
}
{ } {
}
{ } {
}
{ } {
}
}
{}{}{}
gibt, die nicht im Bild der Abbildung liegt, da es keine Funktion auf dem ganzen Kegel ist.





}

\bild{ \begin{center}
\includegraphics[width=5.5cm]{ \bildeinlesungpng {FiberBundle_2} {png} }
\end{center}
\bildtext {Die Fasern einer Abbildung ($M$ ist der Zielbereich, der Definitionsbereich ist die Vereinigung aller Fasern; die Abbildung geht von oben nach unten).} }
\bildlizenz { FiberBundle 2.png } {} {１３２人目'} {ja.wikipedia.org} {CC-BY-SA-3.0} {}






\inputdefinition
{}
{





Zu einem
\definitionsverweis {Morphismus}{}{}
\maabb {{\psi}} { Y } { X
} {}
zwischen
\definitionsverweis {affinen Varietäten}{}{}
bezeichnet man zu einem Punkt

\mavergleichskette
{\vergleichskette
{ P
}
{ \in }{ X
}
{  }{
}
{    }{
}
{   }{
}
}
{}{}{}
das Urbild

\mavergleichskettedisp
{\vergleichskette
{ {\psi}^{-1} (P)
}
{ \subseteq} { Y
}
{ } {
}
{ } {
}
{ } {
}
}
{}{}{}
als die \definitionswort {Faser}{} über $P$. Als abgeschlossene Menge von $Y$ ist sie selbst eine affine Varietät.





}
